% 保留原 Word 的摘要范例。实际论文请用自己的研究摘要替换。
\begin{cnabstract}
随着企业经营环境的不断变化和竞争的日益加剧，企业经营管理水平的提高已成当务之急，人力资源管理作为企业管理中最重要的部分，其水平的高低已成企业发展的关键因素。随着计算机技术特别是互联网技术的发展，管理信息化已成了企业增强核心竞争力的必要手段，开发一个开放、高效、拓展性强的人力资源管理系统对于企业而言具有十分重要的意义。

本文首先分析了人力资源管理以及人力资源管理系统的背景，明确了本文的研究思路\linebreak
和方向，然后，按照软件工程项目的要求，进行了企业的业务需求分析，并通过业务流程\linebreak
整理和问题分析，整理出系统的需求，根据分析结果进行了系统的建模，在需求分析的基\linebreak
础上，进行了系统的设计，包括软件应用架构设计、网络架构设计以及业务逻辑设计，其\linebreak
间，以面向对象的理论为基础，运用\hspace{3.72bp}UML\hspace{3.72bp}进行系统的相关设计，明确了在\hspace{3.72bp}B/S\hspace{3.72bp}模式下采\linebreak
用二层架构来实现目标系统，在系统分析和设计的基础上，进行了系统的数据库设计，包\linebreak
括数据库建模\hspace{0bp plus0.96bp}(E-R\hspace{6bp plus0.96bp}图)\hspace{0bp plus0.96bp}和数据库逻辑设计，最后，对系统进行了开发和运行测试，通过\linebreak
ADO.NET\hspace{4.2bp}和\hspace{4.2bp}AJAX\hspace{4.2bp}等技术的运用，开发了一个用户交互性能良好、功能实用、能满足当\linebreak
前需求以及可扩展的人力资源管理系统，并对开发的系统进行了相关的测试。

在本文的研究中，在系统的架构设计方面以及用户交互体验方面做了较多的工作，目标是开发出一个拓展性强、代码重用度高的系统，以便将来的拓展和维护，同时，改善传统B/S结构用户交互体验差和性能差的状况，提高系统的实用性及易用性，收到了比较好的效果。

\keywords{人力资源管理系统；二层架构；AJAX}
\templatehint{\heiti（注意：关键词是贯穿全文、使用频率比较高的技术名词，3-5个，用分号隔开）}
\end{cnabstract}

\begin{enabstract}
It is great significant to promote the management of the enterprise with not only the\linebreak
changing of environment but also the increasing of competition. However, the management of\linebreak
human resource always plays a crucial part in the enterprise management and the excellent one\linebreak
will lead to a powerful enterprise. Furthermore, as the computer, especially Internet technology\linebreak
have been developed promptly, an open, efficient, and widely expanded human resources\linebreak
management system which has become one of the indispensable manners to enhance the core\linebreak
competitiveness of enterprises should be built urgently.

In this paper, first of all, the background of the human resources management and its system\linebreak
has been analyzed in order to clear out the intention of this research. Secondly, in accordance\linebreak
with the requirements of software engineering project the demands of the enterprise business and\linebreak
its system which are based on the analysis of process and problem in business have been put\linebreak
forward and then the system has been modeled in this chapter. Thirdly, the system including\linebreak
software applications architecture, network architecture and business process is founded on OOA\linebreak
and presented by applying UML and three-layer architecture with B/S in Chapter III. In the\linebreak
fourth chapter, the database involving the modeling (ER diagram) and the logical devise has\linebreak
been schemed. The development employing ADO.NET and AJAX as well as the test of system\linebreak
which is interactive, practical and scalable have been conducted in the fifth chapter.

This study has made great efforts on the architecture design and the interaction to achieve a\linebreak
strongly expanded, highly code-reused, conveniently maintained. This system which is\linebreak
practicable and expedient is so different from the traditional B/S structure based one which is\linebreak
less humanized and capable. In this paper, the main developed parts focus on the management of\linebreak
basic data, personnel alteration, payroll and attendance and systems itself.

\keywordsen{Human resources management system; Three-layer architecture; AJAX}
\end{enabstract}
